\documentclass[11pt,a4paper]{article}
\usepackage[T1]{fontenc}
\usepackage[english]{babel}
\usepackage[margin=26mm,top=23mm,bottom=25mm,footskip=13mm]{geometry}
\usepackage{newtxtext,newtxmath,microtype}
\usepackage[hidelinks]{hyperref}
\hypersetup{pdftitle={Team 4: Peer Review Rebuttal},
  pdfsubject={Introduction to Financial Engineering term project}}
\pagestyle{plain}
\setlength{\parindent}{0pt}
\setlength{\parskip}{4pt}
\raggedbottom

\newenvironment{reviewitem}[1]
  {\par\addvspace{9pt}\noindent\begin{minipage}{\linewidth}%
   \pdfbookmark[2]{Question #1}{question-#1}%
   \textbf{Q#1.}\enspace\ignorespaces}
  {\end{minipage}\par}
\newcommand{\answer}{\par\smallskip\textit{Answer.}\enspace}
\newcommand{\reportref}[1]{\par\smallskip{\small\textit{Report.} #1}}

\begin{document}
\begin{center}
  {\Large Peer Review Rebuttal\par}
  \smallskip
  {\itshape Low-Risk and Coskewness-Based Stock Selection\par}
  \smallskip
  Team 4 \quad Introduction to Financial Engineering
\end{center}

Each month we selected 30 stocks using low beta, low volatility, and negative market coskewness. The baseline portfolio (Capped VW) weights holdings by market capitalization, with a 6.67\% cap per stock.

\pdfbookmark[1]{Peer-review questions}{peer-review}
\section*{Peer-review questions}
Questions follow the original review sheet, with minor typographical corrections. References give the corresponding pages of the final project report.

\begin{reviewitem}{1}
\textit{Incremental Contribution of Coskewness: Although the appendix compares standalone and combined signals, your team noted that the current analysis preserves the original selection results and does not separately identify the effects of all three signals. If your team compares ``low beta + low volatility'' against ``low beta + low volatility + coskewness'' with identical stock counts, sector constraints, and position caps on an unused period, does coskewness still add value?}
\answer With the same 30-stock target, retention rule, nine-stock industry limit, and 6.67\% position cap, adding coskewness raised the full-sample Sharpe ratio from 0.83 to 1.11 (gain 0.281, 95\% confidence interval [0.054, 0.553]). In 2012--2020 the estimated gain was 0.194 (95\% interval $[-0.135,\,0.646]$). Both subsamples were examined retrospectively, leaving performance on unused data untested.
\reportref{Section II-D, p.~2, and Section III-A(4), Table I, Panel B, pp.~3--4.}
\end{reviewitem}

\begin{reviewitem}{2}
\textit{Why Did Optimization Add Little Value? Your team estimated a 36-month covariance matrix for 30 stocks and imposed a common position cap of approximately 6.67\%. Holding the selected stocks fixed while varying the estimation window and the position cap separately, could your team determine whether the limited incremental benefit of optimization results from estimation error, an already defensive stock universe, or restrictive weight limits?}
\answer Under the position cap, performance was largely insensitive to the choice of covariance estimator. Across estimation windows of 24 to 60 months, sample and shrinkage estimators produced Sharpe ratios of 1.095--1.116 under the 6.67\% cap. Without the cap, the 36-month sample estimate yielded a Sharpe ratio of 0.839, and shrinkage improved it by 0.178. Holding the stocks fixed isolates sensitivity to the cap and estimator but leaves the contribution of the defensive screen unidentified.
\reportref{Section IV-A(1), p.~14, and Appendix C-A, Table XI, p.~21.}
\end{reviewitem}

\begin{reviewitem}{3}
\textit{Weaker Later-Period Results and Benchmark Consistency: Your team disclosed the limitations of using the S\&P 500 price index and also compared CRSP and Fama--French total-return benchmarks. If your team consistently uses the S\&P 500 total-return index and additionally control for low-beta and low-volatility exposure in an FF5 framework, how should the excess returns and coskewness premium---particularly during 2012--2020---be interpreted?}
\answer The official S\&P 500 total-return series was unavailable. The CRSP and Fama--French (FF) total-return markets both include dividends and cover broader stock universes. In 2012--2020, Capped VW's Sharpe advantage was 0.136 over CRSP and 0.038 over FF. Neither difference was statistically significant.

Annualized strategy alpha, the excess return unexplained by the factors, fell from 2.25\% under the FF five-factor model (FF5) to $-0.21\%$ after adding betting against beta (BAB). Including both BAB and a low-volatility factor gave 0.01\% alpha for the reconstructed strategy. A separate portfolio that buys high coskewness-score stocks and shorts low-score stocks had FF5-plus-BAB alpha of 3.96\% over the full sample ($t=2.56$) and 2.54\% in 2012--2020 ($t=1.22$). The later-period estimate provides little evidence of a separate coskewness premium.
\reportref{Section III-H, pp.~10--11, Table III, Panel B, p.~15, and Appendix A, p.~20.}
\end{reviewitem}

\begin{reviewitem}{4}
\textit{CC GMVP minimizes volatility, but in the 2007--2009 financial crisis and the 2020 COVID crisis its cumulative losses were larger than those of market-cap weighting. Why do you think losses grew even though volatility was lower?}
\answer The constant-correlation global minimum-variance portfolio (CC GMVP) targets return variance. A smoother return path can still produce a larger cumulative loss if the portfolio participates less in a recovery. In 2007--2009, CC had lower volatility than Capped VW (15.92\% versus 18.42\%) and a shallower drawdown ($-27.03\%$ versus $-29.00\%$), yet lost more over the full episode (23.93\% versus 20.77\%). In February--June 2020, CC again had lower volatility (25.77\% versus 27.52\%) and a larger cumulative loss (15.91\% versus 11.73\%).
\reportref{Section IV-B, p.~14, and Table V, p.~17.}
\end{reviewitem}

\begin{reviewitem}{5}
\textit{Over the full period the low-risk portfolio outperformed the market, but in 2012--2020 its Sharpe advantage was not statistically significant. Given this, can the low-risk selection strategy be considered effective in the long run?}
\answer The evidence supports outperformance over the 215 observed months. The Sharpe difference versus CRSP was 0.423 (95\% interval $[0.128,\,0.762]$), falling to 0.136 in 2012--2020 ($[-0.247,\,0.584]$). The strategy's Sharpe ratio stayed near 1.1 as CRSP's rose from 0.433 to 0.987. The later estimate is consistent with either outperformance or underperformance, so long-run persistence remains uncertain.
\reportref{Section IV-A(3), p.~14, Table III, Panel B, p.~15, and Section V, p.~19.}
\end{reviewitem}

\begin{reviewitem}{6}
\textit{Market-cap weighting had the highest CAGR of the 19 weighting methods. Do you think the main reason the more complex optimizations performed worse is covariance estimation error, or that risk was already well controlled at the stock-selection stage?}
\answer The defensive selection and position cap left little scope for further risk reduction. Across the 19 weighting rules, annualized volatility ranged from 10.71\% to 11.39\%. Changing the covariance estimator had little effect under the cap (Q2). All 18 alternatives trailed Capped VW even before costs, by 0.43--1.54 percentage points in annualized mean return. Capped VW ranked first in compound annual growth rate (CAGR), although none of the net mean-return differences was significant after adjustment for multiple comparisons.
\reportref{Section IV-A(1), pp.~13--14, Table II, p.~13, and Table VI, p.~18.}
\end{reviewitem}

\begin{reviewitem}{7}
\textit{Which selection factor contributed most to performance: low beta, low volatility, or coskewness? Did you test each factor separately?}
\answer We tested all seven signal combinations. Only coskewness significantly predicted next-month returns in the full-sample cross-sectional regression (0.50 percentage points across the full score ranking, $t=2.41$). Adding low beta and low volatility to coskewness reduced volatility from 22.3\% to 11.4\% and the maximum peak-to-trough loss from 64.0\% to 29.0\%. Removing low beta, low volatility, or coskewness lowered Sharpe by 0.256, 0.189, or 0.281, respectively. The pairwise differences between those contributions were statistically inconclusive.
\reportref{Section III-A(4), pp.~3--4, and Table I, Panel A, p.~4.}
\end{reviewitem}

\begin{reviewitem}{8}
\textit{How do you interpret the statistically insignificant Sharpe ratio advantage after 2011? Does this suggest that the strategy's advantage depends on market conditions?}
\answer Relative performance varied across regimes. All 19 allocations lost less than CRSP in the financial crisis and lagged CRSP during the 2020 decline and rebound. The financial crisis accounted for 36.2 of the 50.7 cumulative log-return points gained over CRSP (71\%). Excluding it still left a Sharpe advantage of 0.413 over 194 months (95\% interval $[0.081,\,0.824]$). The weaker post-2011 estimate is consistent with regime dependence.
\reportref{Section IV-B, p.~14, and Table V, p.~17.}
\end{reviewitem}

\begin{reviewitem}{9}
\textit{The S\&P 500 benchmark appears to exclude dividends. Would the conclusions change if you used a total return benchmark that includes reinvested dividends?}
\answer Dividend-inclusive benchmarks reduce the estimated advantage. Against Capped VW's 14.05\% CAGR, the S\&P 500 price index earned 8.61\%, CRSP 10.86\%, and the FF market 11.54\%. The corresponding gaps were 5.44, 3.19, and 2.51 percentage points per year. All 19 allocations exceeded CRSP's CAGR. Differences in index membership also contribute to the change in the estimated gap.
\reportref{Section III-H, pp.~10--11, and Table II, p.~13.}
\end{reviewitem}

\begin{reviewitem}{10}
\textit{In the 2020 crisis (Feb.--Jun.), the portfolios fell about as much as the market at their worst point (\ensuremath{-}19.87\% vs. \ensuremath{-}20.20\%) but ended much lower (\ensuremath{-}11.73\% vs. \ensuremath{-}1.94\%). Why did the low-risk stocks not lose less, and why did they recover more slowly?}
\answer Capped VW absorbed 98\% of the FF market's February--March 2020 decline. Including coskewness in selection had raised the portfolio's pre-crisis estimated market beta from 0.53 to 0.75. With coskewness omitted, downside capture was 71\%. Capped VW then captured only 44\% of the April--June rebound, leaving a five-month return of $-11.73\%$ against $-1.94\%$ for FF despite similar losses at the trough.
\reportref{Section IV-B(a), p.~14, and Table V, p.~17.}
\end{reviewitem}

\begin{reviewitem}{11}
\textit{The minimum-variance portfolio had a lower coskewness score than market-cap weighting (\ensuremath{-}1.06 pp). Is there evidence that coskewness adds value beyond low beta and low volatility?}
\answer Coskewness added value in selection. Adding it to low beta and low volatility raised Sharpe from 0.83 to 1.11 and annualized mean return by 4.12 percentage points ($t=3.18$). The 1.06-percentile-point decrease arose when CC reweighted the same selected stocks. Constraining CC to match Capped VW's coskewness score lowered its Sharpe ratio from 1.069 to 1.048.
\reportref{Section III-A(4), Table I, pp.~3--4, Table VI, p.~18, and Appendix C-B, Table XII, p.~22.}
\end{reviewitem}

\begin{reviewitem}{12}
\textit{Since past covariance estimates may update slowly, would a shorter data window or daily returns help the minimum-variance portfolio in fast crises like 2020?}
\answer Shorter covariance windows worsened crisis performance. With Ledoit--Wolf estimates using 36, 12, and 6 months, February--June 2020 returns were $-10.35\%$, $-12.41\%$, and $-14.63\%$, respectively. Monthly turnover rose from 19.5\% to 24.5\% to 28.9\%. Faster estimates may have reduced exposure after volatility spiked, ahead of the rebound. Daily stock returns were unavailable for this test. Even with daily estimates, portfolio weights could change only at month-end under the existing rule.
\reportref{Appendix C-A, p.~21, discussion following Table XI.}
\end{reviewitem}

\begin{reviewitem}{13}
\textit{Your main benchmark is the S\&P 500 without dividends, but your portfolios include dividends. Against the CRSP market with dividends, your lead shrinks from about 5.4 to 3.2 points a year. Why not use the fair benchmark as the main one?}
\answer We adopted CRSP as the primary benchmark because both CRSP and the strategy include dividends. We retained the FF total-return market as a secondary benchmark and the S\&P 500 price index as a large-cap reference. Against CRSP, Capped VW earned 3.19 percentage points more per year in compound returns, with a higher Sharpe ratio (1.110 versus 0.687) and a shallower maximum drawdown ($-29.00\%$ versus $-51.48\%$).
\reportref{Section III-H, pp.~10--11, Table II, p.~13, and Table V, p.~17.}
\end{reviewitem}

\begin{reviewitem}{14}
\textit{Your lead over the market is significant before 2012 but not after, and you say you can't tell which of your three signals drives the results. Is your strategy mostly the well-known low-risk effect, which has become weaker over time?}
\answer Controlling for BAB reduced later-period FF5 alpha to $-0.21\%$. BAB's own annualized mean return rose from 4.26\% to 9.42\% across subsamples, while the realized coskewness long--short spread fell from 5.96\% to 0.58\%. The matched selection tests identify a full-sample gain from adding coskewness (Q1, Q7). The change in that gain across subsamples was statistically inconclusive.
\reportref{Table I, p.~4, Table III, Panel B, p.~15, and Appendix A, p.~20.}
\end{reviewitem}

\begin{reviewitem}{15}
\textit{During the 2020 Covid crash, your main portfolio lost about 16\% while the S\&P 500 lost about 4\%. Why did a low-risk portfolio lose four times more than the market in a crash, which is exactly when it should protect investors?}
\answer The 16\% figure refers to CC GMVP's February--June return ($-15.91\%$), compared with $-3.88\%$ for the S\&P 500 price index. Their peak-to-trough drawdowns were similar ($-20.12\%$ and $-19.87\%$, respectively). CC's weaker rebound explains the fourfold difference in the five-month loss. Over those same months, the baseline Capped VW portfolio returned $-11.73\%$ and CRSP returned $-3.41\%$.
\reportref{Section IV-B, p.~14, and Table V, p.~17.}
\end{reviewitem}

\begin{reviewitem}{16}
\textit{Several weighting methods were applied to the same stocks, and CC-GMVP didn't clearly improve return or Sharpe ratio. One reason could be that the optimization changes the portfolio's exposure to the scores used to pick the stocks. For example, compared with market-cap weighting it had a higher low-volatility score and a lower coskewness score. If you added a constraint to keep those exposures fixed, what would happen to return and volatility?}
\answer Matching Capped VW's three portfolio-weighted signal scores each month raised CC's CAGR from 13.20\% to 13.29\%. Its annualized mean-return shortfall narrowed from 0.79 to 0.66 percentage points. Volatility rose from 11.08\% to 11.48\%, Sharpe fell from 1.069 to 1.043, and maximum drawdown deepened from $-29.87\%$ to $-32.28\%$.
\reportref{Section IV-A(1), p.~14, and Appendix C-B, Table XII, p.~22.}
\end{reviewitem}

\begin{reviewitem}{17}
\textit{The Sharpe advantage over the market was smaller in 2012--2020 than in 2003--2011. The strategy's Sharpe ratio stayed around 1.1 while the market's went up, but the FF5 alpha also dropped from 9.69\% to 2.25\%. Is this because the portfolio's factor exposures changed, or because the market paid less for the same factors? Did you check this, for example with rolling-window factor regressions?}
\answer We estimated 36- and 60-month rolling FF5 regressions. In the 60-month windows, market beta ranged from 0.51 to 0.80 and annualized alpha fell from 9.20\% in December 2008 to 6.41\% in December 2012 and 3.07\% in December 2020. Allowing all factor exposures to change in January 2012 still yielded a 7.44-percentage-point decline in alpha ($t=-2.79$). Market and investment-factor loadings rose by 0.17 and 0.50, but the joint test of all loading changes was sensitive to lag length ($p=0.006$ to $0.090$). The fall in unexplained return is clear. Its economic cause remains unidentified.
\reportref{Table III, Panel B, p.~15, and Appendix A, equation (40), p.~20.}
\end{reviewitem}

\begin{reviewitem}{18}
\textit{Negative coskewness was used together with low beta and low volatility to select stocks. Depending on how it is defined, negative coskewness can go along with bad returns when the market moves sharply. If you kept the low-beta and low-volatility rules but dropped the coskewness signal, how would downside beta and losses in crash months change?}
\answer Removing coskewness lowered downside beta from 0.79 to 0.66, indicating less sensitivity to the market in months of negative market excess returns. Average portfolio returns improved from $-4.02\%$ to $-3.29\%$ in the 22 bottom-decile market months and from $-9.30\%$ to $-7.10\%$ in the five worst months. The full-sample Sharpe ratio fell from 1.11 to 0.83. The historical return gain from coskewness therefore came with greater downside exposure.
\reportref{Section III-A(5), p.~3, and Table I, p.~4.}
\end{reviewitem}

\begin{reviewitem}{19}
\textit{Slide 44 demonstrates that the Constant Correlation GMVP (CC GMVP) increases the low-volatility tilt (+2.04) while reducing the effective number of holdings ($N_{\mathrm{eff}}$) from 19.8 to 16.9. How does this increased concentration in low-volatility securities influence portfolio diversification during market fluctuations, and does the reduction in $N_{\mathrm{eff}}$ introduce unexpected sector or stock-specific risk?}
\answer CC assigned more weight to a smaller subset of the same 30 stocks. The effective number of holdings ($1/\sum_i w_i^2$) fell from 19.8 to 16.9, and the average weight of the largest industry in each month rose from 26.6\% to 35.2\%. The five largest stocks accounted for 33.3\% under both rules. An industry shock could therefore affect a larger share of the portfolio, although CC's realized volatility was lower in both crises. A 25\% industry cap improved CC's Sharpe ratio by only 0.018 (95\% interval $[-0.054,\,0.099]$). The evidence does not attribute the shortfall to an individual stock or industry.
\reportref{Section IV-C(2), Table VI, p.~18, and Appendix C-B, Table XII, p.~22.}
\end{reviewitem}

\begin{reviewitem}{20}
\textit{Slide 46 compares a 30-stock defensive screen with a 50-stock return-oriented screen (momentum, profitability, issuance). The return-oriented screen increases compound annual growth rate (CAGR) by 2.43 percentage points but also significantly raises volatility (by 7.06 percentage points) and deepens maximum drawdown from \ensuremath{-}29.00\% to \ensuremath{-}48.05\%. Given that the Sharpe ratio declines from 1.110 to 0.858, is it possible to integrate return-oriented factors into a low-risk framework without substantially compromising downside protection?}
\answer We added momentum, distance from the 52-week low, profitability exposure, or net issuance to the defensive screen one at a time, holding the stock count at 30. All four reduced Sharpe (0.913--0.977). The momentum addition slightly reduced maximum drawdown to $-28.45\%$, but CAGR fell to 12.95\%.

An 80\% defensive, 20\% return-oriented blend raised CAGR from 14.05\% to 14.67\% and deepened maximum drawdown from $-29.00\%$ to $-32.22\%$. Sharpe was essentially unchanged in the matched rerun (1.111 to 1.115, 95\% interval for the difference $[-0.102,\,0.123]$). At a 50/50 split, drawdown deepened to $-38.30\%$.
\reportref{Section IV-C(6), p.~19, and Table IX, p.~20.}
\end{reviewitem}

\begin{reviewitem}{21}
\textit{During Crisis I (2007--2009), all defensive strategies substantially limited maximum drawdown (approximately \ensuremath{-}28\% to \ensuremath{-}29\%) compared to the market index (\ensuremath{-}52.56\%). In contrast, during Crisis II (February to June 2020), all 19 strategy allocations underperformed relative to market benchmarks. What structural differences between a prolonged credit crisis (Crisis I) and a rapid, policy-driven rebound (Crisis II) account for the underperformance of low-risk and coskewness strategies during swift market recoveries?}
\answer Capped VW absorbed 58\% of the FF market decline over the 16 months from November 2007 to February 2009 and captured 68\% of the March--May rebound. In 2020 it absorbed 98\% of the two-month decline and captured only 44\% of the April--June rebound. Monthly rebalancing allowed repeated changes during the first decline and only one intervening adjustment during the second. All 19 rules held the same selected stocks each month, limiting what reweighting could recover. The $-52.56\%$ drawdown cited in the question uses the S\&P 500 price index. These two episodes do not isolate the effects of policy intervention or rebalancing frequency.
\reportref{Section IV-B, p.~14, Table V, p.~17, and Figure 10, p.~18.}
\end{reviewitem}

\pdfbookmark[1]{Revisions and replication}{revisions}
\section*{Revisions to the project report}

\textit{Stock selection.} We re-estimated all seven signal combinations. The code in \path{selection/} and inputs in \path{data/selection_inputs/} reproduce both selection files, including 6,450 defensive stock-month records. Some monthly inputs derive from daily data for which the original aggregation code was unavailable (Table I and Appendix E).

\textit{Benchmarks and regimes.} We adopted the CRSP total-return market as the primary benchmark, added BAB and low-volatility controls, estimated rolling factor exposures, and measured decline and rebound capture separately. We also repeated the Sharpe comparison excluding the financial crisis (Tables II, III, and V and Appendix A).

\textit{Portfolio construction.} Beyond the 19 original weighting rules, we evaluated 49 specifications varying estimation windows, position limits, tracking error, signal exposures, and industry concentration. We also tested fourth-signal additions and portfolio blends. Allocation and reporting code is in \path{src/lowrisk_portfolio/} (Tables IX, XI, and XII).

\textit{Data and limitations.} We used project CRSP data, FF factors, and AQR's U.S. BAB series (downloaded November~3, 2026). The cost model assumes 20 basis points per round trip and excludes market impact. Confidence intervals omit adjustment for the original search over signals and portfolio size. The later subsample and the 49 additional specifications reuse historical data. Daily covariance estimation and intramonth rebalancing remain untested (Section V and Appendices C and E).

\end{document}
